\chapter{The Fundamental Group of Manifolds of Negative Curvature}
\label{chap:dc-fundamental-group-negative-curvature}

Let $M^n$ be complete with sectional curvature $K<0$. By the Hadamard theorem,
its universal cover is diffeomorphic to $\mathbb{R}^n$; consequently
$\pi_k(M)=0$ for $k\geq 2$. The results below constrain $\pi_1(M)$ when $M$ is
compact.

\section{Existence of closed geodesics}

Let $\pi:\tilde M\to M$ be the universal covering with the covering metric, and
let $A(\tilde M)$ be its group of covering transformations. Every nonidentity
element of $A(\tilde M)$ is a fixed-point-free isometry, and a choice of
$\tilde p\in\pi^{-1}(p)$ determines an isomorphism
$\pi_1(M;p)\cong A(\tilde M)$. Explicitly, for $g\in\pi_1(M;p)$ and
$\tilde q\in\tilde M$, choose a path $\tilde\sigma$ from $\tilde q$ to
$\tilde p$, set $\sigma=\pi\circ\tilde\sigma$ and
$\beta=\sigma g\sigma^{-1}$, and define
\[
  \alpha_{\tilde p}(\tilde q)=\tilde\beta_{\tilde q}(1),
\]
where $\tilde\beta_{\tilde q}$ is the lift of $\beta$ starting at $\tilde q$.

\begin{definition}[Free homotopy class]
  \label{def:dc-ch12-2-1}
  \dcref{ch12:2.1}
  Closed paths $f,g:I\to M$ are \emph{freely homotopic} if there is a map
  $F:I\times I\to M$ such that
  \[
    F(0,t)=f(t),\qquad F(1,t)=g(t),\qquad F(s,0)=F(s,1).
  \]
  The corresponding equivalence classes form the set $C_1(M)$. Unlike based
  homotopy classes, free homotopy classes allow the base point to vary.

  A \emph{closed geodesic} is a closed curve geodesic at every point. A
  \emph{geodesic lasso} is closed and geodesic except at one nonregular point.
\end{definition}

\begin{theorem}[Cartan's closed-geodesic theorem]
  \label{thm:dc-ch12-2-2}
  \dcref{ch12:2.2}
  If $M$ is compact and $\mathcal L\in C_1(M)$ is nonconstant, then
  $\mathcal L$ contains a closed geodesic.
\end{theorem}
\begin{proof}
  \sketch
  Let $d$ be the infimum of the lengths of piecewise differentiable curves in
  $\mathcal L$; nontriviality gives $d>0$. Choose broken geodesics
  $\gamma_j:[0,1]\to M$, parametrized proportionally to arc length, with
  $\ell(\gamma_j)\to d$. If $L=\sup_j\ell(\gamma_j)$, then
  \[
    d(\gamma_j(t_1),\gamma_j(t_2))
    \leq \int_{t_1}^{t_2}|\gamma_j'(t)|\,dt
    \leq L(t_2-t_1).
  \]
  Compactness and equicontinuity give a uniformly convergent subsequence with
  closed limit $\gamma_0$.

  Choose $0=t_0<\cdots<t_k=1$ so that each corresponding arc of $\gamma_0$
  lies in a totally normal neighborhood, and join successive points
  $\gamma_0(t_i)$ by the unique minimizing segments to obtain
  $\gamma\in\mathcal L$. If $\ell(\gamma)>d$, set
  $\varepsilon=(\ell(\gamma)-d)/(2k+1)$. For large $j$,
  \[
    \ell(\gamma_j)-d<\varepsilon,
    \qquad d(\gamma_j(t),\gamma_0(t))<\varepsilon,
  \]
  so
  \[
    \sum_{i=1}^k\bigl(\ell(\gamma_j^i)+2\varepsilon\bigr)
    <\ell(\gamma)=\sum_{i=1}^k\ell(\gamma^i).
  \]
  Some minimizing segment $\gamma^i$ would then be longer than a competitor,
  a contradiction. Thus $\ell(\gamma)=d$. After arc-length parametrization,
  $\gamma$ has no corner: a corner inside a convex ball could be shortened by
  replacing the two incident arcs with the minimizing chord, without changing
  the free homotopy class. Hence $\gamma$ is a closed geodesic.
\end{proof}

\begin{remark}[Failure without compactness]
  \label{rem:dc-ch12-2-3}
  \dcref{ch12:2.3}
  Compactness is necessary.
\end{remark}
\begin{proof}

  \sketch
  A surface of revolution generated by a curve
  asymptotic to its axis has nontrivial free homotopy classes containing curves
  of arbitrarily small length, so those classes have no length minimizer.
\end{proof}

\begin{remark}[Simply connected compact case]
  \label{rem:dc-ch12-2-4}
  \dcref{ch12:2.4}
  A compact simply connected Riemannian manifold also has a closed geodesic,
  although this does not follow from a nonconstant free homotopy class.
\end{remark}

\begin{definition}[Translation along a geodesic]
  \label{def:dc-ch12-2-5}
  \dcref{ch12:2.5}
  A fixed-point-free isometry $f:\tilde M\to\tilde M$ is a
  \emph{translation} if some geodesic $\tilde\gamma$ is invariant under $f$,
  meaning
  \[
    f\bigl(\tilde\gamma(\mathbb R)\bigr)=\tilde\gamma(\mathbb R).
  \]
  In that case $f$ is a \emph{translation along} $\tilde\gamma$.
\end{definition}

\begin{proposition}[Covering transformations are translations]
  \label{prop:dc-ch12-2-6}
  \dcref{ch12:2.6}
  Let $M$ be compact. Every nonidentity covering transformation of
  $\tilde M$ is a translation.
\end{proposition}
\begin{proof}
  \sketch
  Let $\alpha$ correspond to $g\in\pi_1(M;p)$ after choosing
  $\tilde p\in\pi^{-1}(p)$. By \cref{thm:dc-ch12-2-2}, the associated free
  homotopy class contains a closed geodesic $\gamma$. If $q\in\gamma$ and
  $\sigma$ joins $p$ to $q$, lift $\sigma$ from $\tilde p$ to an endpoint
  $\tilde q$, and lift $\gamma$ from $\tilde q$. The deck transformation
  $\alpha_{\tilde q}$ associated with $[\gamma]\in\pi_1(M;q)$ satisfies
  \[
    \alpha(\tilde q)=\alpha_{\tilde q}(\tilde q).
  \]
  Thus $\alpha\alpha_{\tilde q}^{-1}$ has a fixed point and is the identity, so
  $\alpha=\alpha_{\tilde q}$. Uniqueness of path lifting then yields
  \[
    \alpha(\tilde\gamma(s))
    =\alpha_{\tilde q}(\tilde\gamma(s))\in\tilde\gamma(\mathbb R),
  \]
  proving invariance.
\end{proof}

\section{Preissman's Theorem}

A \emph{geodesic triangle} consists of three unit-speed minimizing geodesic
segments
\[
  \gamma_1:[0,\ell_1]\to M,\qquad
  \gamma_2:[0,\ell_2]\to M,\qquad
  \gamma_3:[0,\ell_3]\to M,
\]
with $\gamma_i(\ell_i)=\gamma_{i+1}(0)$ for $i=1,2$ and
$\gamma_3(\ell_3)=\gamma_1(0)$. Its interior angles are
\[
  \sphericalangle(-\gamma_i'(\ell_i),\gamma_{i+1}'(0))\quad(i=1,2),
  \qquad
  \sphericalangle(-\gamma_3'(\ell_3),\gamma_1'(0)).
\]

\begin{lemma}[Comparison for geodesic triangles when $K\leq 0$]
  \label{lem:dc-ch12-3-1}
  \uses{prop:dc-ch10-2-5}
  \dcref{ch12:3.1}
  Let $\tilde M$ be complete, simply connected, and satisfy $K\leq0$. Points
  $a,b,c\in\tilde M$ determine a unique geodesic triangle. Let
  $\alpha,\beta,\gamma$ be the angles at $a,b,c$, and let $A,B,C$ be the
  lengths of the opposite sides. Then
  \begin{enumerate}
    \item $A^2+B^2-2AB\cos\gamma\leq C^2$;
    \item $\alpha+\beta+\gamma\leq\pi$.
  \end{enumerate}
  Both inequalities are strict when $K<0$.
\end{lemma}
\begin{proof}
  \sketch
  Let $\gamma_A,\gamma_B,\gamma_C$ be the sides of lengths $A,B,C$, with the
  first two issuing from $c$, and set
  $\Gamma_i=\exp_c^{-1}\circ\gamma_i$ for $i=A,B,C$. Let $\Gamma_0$ be the
  straight segment joining the endpoints of $\Gamma_C$. The radial sides satisfy
  \[
    A=\ell(\Gamma_A),\qquad B=\ell(\Gamma_B),
  \]
  while Euclidean geometry in $T_c\tilde M$ gives
  \[
    \ell(\Gamma_0)^2=A^2+B^2-2AB\cos\gamma.
  \]
  Since $\ell(\Gamma_0)\leq\ell(\Gamma_C)$, Rauch comparison
  (\cref{prop:dc-ch10-2-5}) gives
  \[
    A^2+B^2-2AB\cos\gamma
    \leq\ell(\Gamma_C)^2
    \leq\ell(\gamma_C)^2=C^2,
  \]
  with strict final inequality when $K<0$. This proves the first assertion.

  The side lengths satisfy the triangle inequalities, so they determine a
  Euclidean comparison triangle with opposite angles
  $\alpha',\beta',\gamma'$. Applying the first assertion at each vertex yields
  $\alpha\leq\alpha'$, $\beta\leq\beta'$, and
  $\gamma\leq\gamma'$, all strict when $K<0$. Since
  $\alpha'+\beta'+\gamma'=\pi$, the angle-sum assertion follows.
\end{proof}

For the remainder of the chapter, $M$ is complete with $K<0$, and
$\pi:\tilde M\to M$ is its universal covering with the covering metric.

\begin{theorem}[Preissman's cyclic-subgroup theorem]
  \label{thm:dc-ch12-3-2}
  \dcref{ch12:3.2}
  If $M$ is compact, every nontrivial abelian subgroup of $\pi_1(M)$ is
  infinite cyclic.
\end{theorem}
\begin{proof}
  \sketch
  Choose a nonidentity element of the abelian subgroup $H$. By
  \cref{prop:dc-ch12-2-6}, it translates along a geodesic
  $\tilde\gamma$. The uniqueness in \cref{lem:dc-ch12-3-3} and commutativity,
  through \cref{lem:dc-ch12-3-4}, imply that every element of $H$ preserves
  $\tilde\gamma$. Now \cref{lem:dc-ch12-3-5} shows that $H$ is infinite
  cyclic.
\end{proof}

\begin{lemma}[Uniqueness of an invariant geodesic]
  \label{lem:dc-ch12-3-3}
  \uses{lem:dc-ch12-3-1}
  \dcref{ch12:3.3}
  If $K<0$ and a nonidentity translation $f:\tilde M\to\tilde M$ preserves
  $\tilde\gamma$, then $\tilde\gamma$ is the unique geodesic preserved by $f$.
\end{lemma}
\begin{proof}
  \sketch
  Suppose $f$ preserves distinct geodesics $\tilde\gamma_1$ and
  $\tilde\gamma_2$. They are disjoint: if their intersection were a single
  point, it would be fixed by $f$, while two intersection points would force
  the geodesics to coincide. Join
  $\tilde p_i\in\tilde\gamma_i$ by a minimizing geodesic. The resulting
  quadrilateral with vertices
  $\tilde p_1,f(\tilde p_1),f(\tilde p_2),\tilde p_2$ has interior angles
  $\alpha,\pi-\alpha',\beta,\pi-\beta'$. Since $f$ is an isometry,
  $\alpha=\alpha'$ and $\beta=\beta'$, so their sum is $2\pi$. Dividing the
  quadrilateral into two triangles forces one triangle to have angle sum at
  least $\pi$, contradicting the strict case of
  \cref{lem:dc-ch12-3-1}.
\end{proof}

\begin{lemma}[A commuting isometry preserves the translation axis]
  \label{lem:dc-ch12-3-4}
  \dcref{ch12:3.4}
  If $K<0$, $f$ is a nonidentity translation along $\tilde\gamma$, and the
  fixed-point-free isometry $g$ commutes with $f$, then $g$ is a translation
  along $\tilde\gamma$.
\end{lemma}
\begin{proof}
  \sketch
  Commutativity gives
  \[
    f\bigl(g(\tilde\gamma)\bigr)
    =g\bigl(f(\tilde\gamma)\bigr)=g(\tilde\gamma).
  \]
  Thus $g(\tilde\gamma)$ is another $f$-invariant geodesic, so
  \cref{lem:dc-ch12-3-3} gives $g(\tilde\gamma)=\tilde\gamma$.
\end{proof}

\begin{lemma}[A subgroup preserving one geodesic is infinite cyclic]
  \label{lem:dc-ch12-3-5}
  \uses{prop:dc-ch12-2-6}
  \dcref{ch12:3.5}
  If every element of a nontrivial subgroup $H\subset\pi_1(M)$, acting on
  $\tilde M$, preserves a fixed geodesic $\tilde\gamma$, then $H$ is infinite
  cyclic.
\end{lemma}
\begin{proof}
  \sketch
  Orient $\tilde\gamma$, fix $\tilde p\in\tilde\gamma$, and assign to
  $h\in H$ the signed displacement
  \[
    \theta(h)=\pm d(\tilde p,h(\tilde p)).
  \]
  Because every $h$ acts isometrically on $\tilde\gamma$, this defines a
  homomorphism $\theta:H\to(\mathbb R,+)$. It is injective: equality of two
  images would make $h_1h_2^{-1}$ fix $\tilde p$. Hence $H$ is isomorphic to
  an additive subgroup of $\mathbb R$. Such a subgroup is either dense or
  infinite cyclic; proper discontinuity of the deck action excludes density.
\end{proof}

\begin{example}[Surface of genus two times a circle]
  \label{ex:dc-ch12-3-6}
  \dcref{ch12:3.6}
  Let $N$ be a surface of genus two. Then $N\times S^1$ admits no negatively
  curved metric.
\end{example}
\begin{proof}

  \sketch
  For any infinite cyclic subgroup $C\subset\pi_1(N)$, the subgroup
  $C\oplus\mathbb Z\subset\pi_1(N\times S^1)$ is abelian but not cyclic,
  contradicting \cref{thm:dc-ch12-3-2}.
\end{proof}

\begin{example}[The $m$-torus]
  \label{ex:dc-ch12-3-7}
  \dcref{ch12:3.7}
  For $m\geq2$, the torus $T^m=(S^1)^m$ admits no negatively curved metric.
\end{example}
\begin{proof}

  \sketch
  Its fundamental group is $\mathbb Z^m$, a noncyclic abelian group, contrary
  to \cref{thm:dc-ch12-3-2}.
\end{proof}

\begin{theorem}[Nonabelian fundamental group]
  \label{thm:dc-ch12-3-8}
  \dcref{ch12:3.8}
  If $M$ is compact and $K<0$, then $\pi_1(M)$ is not abelian.
\end{theorem}
\begin{proof}
  \sketch
  Suppose $\pi_1(M)$ were abelian. If it were trivial, every geodesic in
  $\tilde M$ would be invariant under all deck transformations, contradicting
  \cref{lem:dc-ch12-3-9}. Otherwise, choose a nonidentity deck transformation
  and its unique invariant geodesic. By \cref{lem:dc-ch12-3-4}, every deck
  transformation preserves that geodesic, again contradicting
  \cref{lem:dc-ch12-3-9} and compactness.
\end{proof}

\begin{lemma}[An invariant geodesic forces noncompactness]
  \label{lem:dc-ch12-3-9}
  \uses{lem:dc-ch12-3-1}
  \dcref{ch12:3.9}
  If $M$ is complete, $K\leq0$, and some geodesic in $\tilde M$ is invariant
  under every element of $A(\tilde M)$, then $M$ is noncompact.
\end{lemma}
\begin{proof}
  \sketch
  Let $\tilde\gamma$ be invariant, fix $\tilde p\in\tilde\gamma$, and for
  $t>0$ let $\tilde\beta:[0,t]\to\tilde M$ be the unit-speed geodesic from
  $\tilde p$ perpendicular to $\tilde\gamma$. Write
  $\beta=\pi\circ\tilde\beta$, $p=\pi(\tilde p)$, and let $\alpha_t$ be a
  minimizing geodesic from $\beta(t)$ to $p$. The lift $\tilde\alpha_t$ from
  $\tilde\beta(t)$ ends on $\tilde\gamma$ because the latter is invariant under
  every deck transformation. The comparison inequality
  \cref{lem:dc-ch12-3-1} gives
  $\ell(\tilde\alpha_t)\geq\ell(\tilde\beta)=t$, whereas
  \[
    \ell(\tilde\alpha_t)=\ell(\alpha_t)
    \leq\ell(\beta)=\ell(\tilde\beta)=t.
  \]
  Thus $d(\beta(t),p)=t$ for every $t>0$, so $M$ is unbounded and hence
  noncompact.
\end{proof}

\begin{theorem}[Byers's solvable-subgroup theorem]
  \label{thm:dc-ch12-3-10}
  \uses{thm:dc-ch12-3-8, lem:dc-ch12-3-5}
  \dcref{ch12:3.10}
  If $M$ is compact, $K<0$, and $H\neq\{e\}$ is a solvable subgroup of
  $\pi_1(M)$, then $H$ is infinite cyclic. Moreover, $\pi_1(M)$ has no cyclic
  subgroup of finite index.
\end{theorem}
\begin{proof}
  \sketch
  Choose a shortest solvable series
  \[
    H=H_0\supset H_1\supset\cdots\supset H_{k-1}\supset H_k=\{e\},
  \]
  where $H_{i+1}\triangleleft H_i$ and $H_i/H_{i+1}$ is abelian. Then
  $H_{k-1}$ is nontrivial and abelian, hence infinite cyclic by
  \cref{thm:dc-ch12-3-2}. Let $g$ generate it and let $\tilde\gamma$ be the
  $g$-invariant geodesic. If $k=1$, this already proves the first assertion.
  Assume $k\geq2$. For $a\in H_{k-2}$ and nonidentity
  $b\in H_{k-1}$, normality gives an integer $n$ such that
  \[
    a^{-1}b^{-1}ab=g^n.
  \]
  Since both $b$ and $g^n$ preserve $\tilde\gamma$, $b^{-1}$ preserves
  $a(\tilde\gamma)$. Uniqueness of the invariant geodesic gives
  $a(\tilde\gamma)=\tilde\gamma$. Hence every element of $H_{k-2}$ preserves
  $\tilde\gamma$, and \cref{lem:dc-ch12-3-5} makes $H_{k-2}$ infinite
  cyclic. Repeating upward through the series proves that $H$ is infinite
  cyclic.

  Suppose now that a cyclic subgroup $H\subset\pi_1(M)$ has finite index. If
  $H=\{e\}$, then $\pi_1(M)$ is finite; \cref{thm:dc-ch12-3-2} forces it to
  be trivial, contrary to \cref{thm:dc-ch12-3-8}. Thus
  $H=\langle g\rangle$ with $g\neq e$. It is proper, since otherwise
  $\pi_1(M)$ would be abelian. If $\tilde\gamma$ is the $g$-invariant geodesic
  and $a\notin H$, finite index gives $a^n\in H$ for some $n>0$.
  Preissman's theorem applied to $\langle a\rangle$ implies $a^n\neq e$, so
  $a^n=g^m$ for some $m\neq0$. Both $\tilde\gamma$ and
  $a(\tilde\gamma)$ are invariant under $a^n$, and uniqueness yields
  $a(\tilde\gamma)=\tilde\gamma$. Thus all of $\pi_1(M)$ preserves
  $\tilde\gamma$, so \cref{lem:dc-ch12-3-5} would make $\pi_1(M)$ infinite
  cyclic, contradicting \cref{thm:dc-ch12-3-8}.
\end{proof}

\begin{remark}[Further nonpositive-curvature structure]
  \label{rem:dc-ch12-3-11}
  \dcref{ch12:3.11}
  The results above use only part of the structure relating the geometry of a
  nonpositively curved manifold to its fundamental group.
\end{remark}
